\emph{Step 1 (the two repairs).} For $A = \Cl(1,3)$, $p - q = -2 \equiv 2 \pmod 4$, so by Theorem \claimref{rule2.where-obstruction-sits-commuting-square}(a) the spacelike extension $\Cl(1,4)$ is obstructed, and by (c) and (d) its Rule 2 repairs are $\Cl(2,4)$ (adjoined $\nu$ of the opposite sign, $\nu^{2} = +1$, the adjoined plane hyperbolic) and $\Cl(1,5)$ ($\nu^{2} = -1$, the adjoined plane definite). For the mirror $\Cl(3,1)$, $p - q = 2$, the obstructed extension is again $\Cl(p,q{+}1) = \Cl(3,2)$ (in this convention its new generator squares like time), with repairs $\Cl(4,2)$ ($\nu^{2} = +1$, opposite sign, the Rule 1 ascent) and $\Cl(3,3)$ ($\nu^{2} = -1$, same sign). The check recomputes these four signatures from the rule of (a).

\emph{Step 2 (the Cartan decomposition and its centre).} In $\Cl(p,q)$ the bivectors $e_{i}e_{j}$, $i < j$, span a Lie algebra under the commutator isomorphic to $\so(p,q)$ (\claimref{std-algebra.bivector-lie-algebra}). Since $(e_{i}e_{j})^{2} = -e_{i}^{2}e_{j}^{2}$, a bivector squares to $-1$ when its two factors square alike and to $+1$ otherwise; the first kind span $\mathfrak{k} \cong \so(p) \oplus \so(q)$, the maximal compact subalgebra (for $(2,4)$ of dimension $7$, matching the $SO(2) \times SO(4)$ of Fact \claimref{standard-model.imported-maximal-compact-subgroups}), and the second span $\mathfrak{p}$. The centre of $\mathfrak{k}$ is the kernel of the linear map $X \mapsto ([X, k])_{k \in \mathfrak{k}}$ on $\mathfrak{k}$; the check computes its rank over $\Q$ and finds kernel dimension $1$ for $(2,4)$ and $(4,2)$ and $0$ for $(1,5)$ and $(3,3)$. By hand: $\so(p) \oplus \so(q)$ has centre $\so(2)$ for each factor with index $2$ and nothing from factors with index $1$ or at least $3$, which gives the same four values. For $(2,4)$ the central element is $Z = e_{0}\nu$, the product of the two generators squaring $+1$, that is of the spacetime time direction and the repair generator; it commutes with every rotation, and the check verifies $Z^{2} = -1$ and $[Z,[Z,x]] = -4x$ for every boost $x$, so $J = \tfrac12\mathrm{ad}(Z)$ satisfies $J^{2} = -1$ on $\mathfrak{p}$ (the Hermitian-type structure). In $\Cl(4,2)$ the two generators squaring $-1$ are the base time direction and the extension generator, their product is the central rotation, and $e_{\mathrm{time}}\nu$ is a boost, $(e_{\mathrm{time}}\nu)^{2} = +1$ (also checked); so the identification of the central plane as the plane of time and $\nu$ holds in the $\Cl(1,3)$ convention only, while in both conventions it is the plane of time and the new generator squaring like time. In $\Cl(1,5)$ the plane of $e_{0}$ and $\nu$ has $(e_{0}\nu)^{2} = -e_{0}^{2}\nu^{2} = +1$: it is a boost plane, and the check confirms that every rotation fails to commute with some other rotation, so no rotation is central.

\emph{Step 3 (existence for the opposite-sign repair).} Proposition \claimref{field-theory.twistor-attachment} constructs, on the conformal spinor space with its $\mathfrak{u}(2,2)$ action, positive-energy modules: the helicity-$s$ module has energies $|s| + 1 + n$, bounded below by $|s| + 1$, so $\mathrm{Spin}^{+}(2,4) \cong SU(2,2)$ has representations of positive energy. That the conformal energy is the compact generator spanning the $\so(2)$ centre of $\mathfrak{k}$ in Step 2 (the conformal Hamiltonian $\tfrac12(P_{0} + K_{0})$ in the usual notation) is the standard identification; the store does not state it, and it is argued here, not proved. Existence does not depend on it.

\emph{Step 4 (non-existence for the same-sign repair, given the import).} Let $\pi$ be a representation of positive energy of $G = \mathrm{Spin}(p,q)$ with $\mathfrak{g} = \so(p,q)$ simple and noncompact, and let $C_{\pi} = \{X \in \mathfrak{g} : -I_{4}\,d\pi(X) \text{ is bounded below}\}$ (scalars in $\mathrm{span}\{1, I_{4}\}$). It is invariant under $\mathrm{Ad}(G)$ (conjugation by $\pi(g)$ preserves semiboundedness) and is a convex cone (a sum of two operators bounded below is bounded below on the common domain of smooth vectors). It is not $0$, by hypothesis. It is not all of $\mathfrak{g}$. Otherwise $X$ and $-X$ both lie in it for every $X$, so every $d\pi(X)$ is bounded. Since $\mathfrak{g}$ has real rank at least one, it contains $H, E$ with $[H, E] = 2E$ and $E \neq 0$; unitarity gives $\|d\pi(E)\| = \|\pi(e^{tH})d\pi(E)\pi(e^{-tH})\| = \|d\pi(\mathrm{Ad}(e^{tH})E)\| = e^{2t}\|d\pi(E)\|$ for every real $t$, so $d\pi(E) = 0$. The kernel of $d\pi$ is then a nonzero ideal of the simple $\mathfrak{g}$, hence all of $\mathfrak{g}$, and $\pi$ is trivial on the identity component, against the definition. Its closure is a closed invariant convex cone, and it is not $\mathfrak{g}$ either, since a convex set in a finite-dimensional space has the same interior as its closure, so a dense $C_{\pi}$ would be all of $\mathfrak{g}$. The imported classification (Fact \claimref{synthesis.imported-invariant-cone-classification}) states that a simple real Lie algebra contains a closed invariant convex cone other than $0$ and $\mathfrak{g}$ only if it is of Hermitian type, that is only if the centre of its maximal compact subalgebra is nonzero. By Step 2, $\so(1,5)$ and $\so(3,3)$ are not of Hermitian type. Hence $\mathrm{Spin}(1,5)$ and $\mathrm{Spin}(3,3)$ have no representation of positive energy. (The requirement of nontriviality on the identity component is needed: these groups are disconnected, and a representation through the component group alone has $d\pi = 0$.)

\emph{Step 5 (conclusion and scope).} Imposing the spectrum condition on the symmetry group of the completed algebra therefore admits $\nu^{2} = +1$ (the Rule 1 ascent) and excludes $\nu^{2} = -1$ for the spacetime algebra and for its mirror, given the import of Step 4. The principle is a requirement on the dynamics (its generator is bounded below), not a statement about the fermion representations; it is not derived here from equations of motion, and the series must still take the spectrum condition as given, as its twistor attachment already does. The base Lorentz algebra $\so(1,3)$ has zero centre in $\mathfrak{k}$ (also checked), so the discrimination appears only after the completion. Labels: Steps 1 and 2 are proved and checked exactly; Step 3 is the store's proved construction; Step 3's identification of the energy generator is argued (standard); Step 4 is proved from the one named import, Fact \claimref{synthesis.imported-invariant-cone-classification}.
